\documentclass[a4paper,11pt]{article}

\usepackage[utf8]{inputenc}
\usepackage[T1]{fontenc}
\usepackage[ngerman]{babel}
\usepackage{amsmath,amssymb}
\usepackage{geometry}
\geometry{margin=2.0cm}
\usepackage{enumitem}
\usepackage{fancyhdr}
\usepackage{xcolor}
\usepackage{tikz}
\usepackage{titlesec}

%================= Dark-Mode Farben (Catppuccin Mocha) =================
\definecolor{base}{HTML}{1E1E2E}     % Hintergrund
\definecolor{fg}{HTML}{CDD6F4}       % Text (heller Vordergrund)
\definecolor{accent}{HTML}{89B4FA}   % Überschriften / Akzent
\definecolor{accent2}{HTML}{F5C2E7}  % Schwierigkeitsmarkierung
\definecolor{gridcolor}{HTML}{45475A}% Karo-Gitter (dezent)

\pagecolor{base}

\newcommand{\transp}{^{\mathsf{T}}}
\newcommand{\norm}[1]{\left\lVert #1 \right\rVert}
\newcommand{\inner}[2]{\left\langle #1,\, #2 \right\rangle}
\newcommand{\rang}{\operatorname{rang}}
\newcommand{\spur}{\operatorname{spur}}
\newcommand{\diag}{\operatorname{diag}}

% Schwierigkeitsmarkierung
\newcommand{\schwierig}[1]{\hfill\textnormal{\color{accent2}\itshape (#1)}}

% Karierte Antwortfläche (Karopapier) -- #1 = Höhe des Bereichs (z. B. 5cm).
% Durchgehendes Karo-Raster, keine einzelnen Schreiblinien, kein \hrulefill.
\newcommand{\karo}[1]{\par\vspace{0.3cm}\noindent
  \begin{tikzpicture}
    \draw[step=5mm, gridcolor, very thin] (0,0) grid (\linewidth,#1);
  \end{tikzpicture}\par\vspace{0.3cm}}

%================= Überschriften in Akzentfarbe =================
\titleformat{\section}{\color{accent}\normalfont\Large\bfseries}{}{0pt}{}

%================= Kopf-/Fußzeile =================
\pagestyle{fancy}
\fancyhf{}
\lhead{\color{fg}\small Fortgeschrittene Lineare Algebra und Analytische Geometrie}
\rhead{\color{fg}\small Übungsblatt 01}
\cfoot{\color{fg}\thepage}
\renewcommand{\headrule}{{\color{gridcolor}\hrule height 0.4pt}}

\setlength{\parindent}{0pt}

\begin{document}
\color{fg}

\thispagestyle{fancy}

\begin{center}
  {\color{accent}\LARGE\bfseries Übungsaufgaben: Matrizen, Zerlegungen \& Geometrie}\\[0.4em]
  {\large Fortgeschrittene Lineare Algebra und Analytische Geometrie}\\[0.8em]
\end{center}

\noindent
\begin{tabular}{@{}p{0.6\textwidth}p{0.35\textwidth}@{}}
  \textbf{Name:} \textcolor{fg}{\hrulefill} & \textbf{Datum:} \textcolor{fg}{\hrulefill}
\end{tabular}

\vspace{0.4em}
{\color{gridcolor}\hrule}
\vspace{0.6em}

\textit{Hinweise: Schwierigkeitsgrad ist je Aufgabe in Klammern angegeben
(\emph{leicht}, \emph{mittel}, \emph{schwer}). Rechenwege bitte nachvollziehbar
notieren. Notation wie in der Vorlesung: $E$ bezeichnet die Einheitsmatrix,
$\norm{\cdot}$ die euklidische ($L_2$-)Norm, sofern nicht anders angegeben.
Lösungen frei auf dem Karo-Gitter notieren.}

\vspace{0.8em}

%==================================================================
\section*{Aufgabe 1 -- Quickies (Wahr oder Falsch)\schwierig{leicht}}
Entscheiden Sie für jede Aussage, ob sie \emph{wahr} oder \emph{falsch} ist,
und begründen Sie kurz.
\begin{enumerate}[label=\alph*)]
  \item Jede orthogonale Matrix $Q$ erfüllt $Q^{-1} = Q\transp$.
  \item Die Singulärwerte einer Matrix können negativ sein.
  \item Eine symmetrische positiv definite Matrix besitzt ausschließlich positive Eigenwerte.
  \item Für \emph{jede} quadratische Matrix existiert eine Eigenwertzerlegung $A = S\Lambda S^{-1}$.
  \item Die Spur einer quadratischen Matrix ist gleich der Summe ihrer Eigenwerte.
  \item Eine Hyperebene im $\mathbb{R}^n$ ist stets ein Untervektorraum.
\end{enumerate}
\karo{5.5cm}

%==================================================================
\section*{Aufgabe 2 -- Normen\schwierig{leicht}}
Gegeben sei der Vektor
\[
  \vec{x} = \begin{pmatrix} 3 \\ -4 \\ 12 \end{pmatrix} \in \mathbb{R}^3 .
\]
\begin{enumerate}[label=\alph*)]
  \item Berechnen Sie die $L_1$-Norm $\norm{\vec{x}}_1$.
  \item Berechnen Sie die $L_2$-Norm $\norm{\vec{x}}_2$.
  \item Berechnen Sie die Maximum-Norm $\norm{\vec{x}}_\infty$.
\end{enumerate}
\karo{5cm}

%==================================================================
\section*{Aufgabe 3 -- Inneres Produkt und Orthogonalität\schwierig{mittel}}
Gegeben seien
\[
  \vec{x} = \begin{pmatrix} 1 \\ 2 \\ 2 \end{pmatrix},\qquad
  \vec{y} = \begin{pmatrix} 2 \\ -2 \\ 1 \end{pmatrix}.
\]
\begin{enumerate}[label=\alph*)]
  \item Berechnen Sie das Skalarprodukt $\inner{\vec{x}}{\vec{y}} = \vec{x}\transp\vec{y}$.
        Sind $\vec{x}$ und $\vec{y}$ orthogonal?
  \item Bestimmen Sie $\norm{\vec{x}}$ und $\norm{\vec{y}}$ und überprüfen Sie die
        Cauchy-Schwarz-Ungleichung $|\inner{\vec{x}}{\vec{y}}| \le \norm{\vec{x}}\cdot\norm{\vec{y}}$.
  \item Verifizieren Sie den Satz des Pythagoras $\norm{\vec{x}+\vec{y}}^2 = \norm{\vec{x}}^2 + \norm{\vec{y}}^2$.
  \item Bestimmen Sie den Winkel $\alpha$ zwischen
        $\vec{u} = (1,1,0)\transp$ und $\vec{v} = (1,0,1)\transp$.
\end{enumerate}
\karo{6.5cm}

%==================================================================
\section*{Aufgabe 4 -- CR-Zerlegung\schwierig{mittel}}
Gegeben sei
\[
  A = \begin{pmatrix} 1 & 2 & 3 \\ 2 & 5 & 7 \\ 3 & 7 & 10 \end{pmatrix}.
\]
\begin{enumerate}[label=\alph*)]
  \item Bestimmen Sie den Rang von $A$.
  \item Geben Sie eine Basis des Spaltenraums an.
  \item Bestimmen Sie die CR-Zerlegung $A = CR$.
\end{enumerate}
\karo{6.5cm}

%==================================================================
\section*{Aufgabe 5 -- LU-Zerlegung\schwierig{mittel}}
Gegeben sei
\[
  A = \begin{pmatrix} 2 & 1 & 1 \\ 4 & 3 & 3 \\ 8 & 7 & 9 \end{pmatrix}.
\]
\begin{enumerate}[label=\alph*)]
  \item Bestimmen Sie eine LU-Zerlegung $A = LU$.
  \item Bestimmen Sie $\det(A)$.
  \item Berechnen Sie $A^{-1}$.
\end{enumerate}
\karo{8cm}

%==================================================================
\section*{Aufgabe 6 -- Cholesky-Zerlegung\schwierig{mittel}}
Gegeben sei
\[
  A = \begin{pmatrix} 4 & 2 & 2 \\ 2 & 5 & 5 \\ 2 & 5 & 14 \end{pmatrix}.
\]
\begin{enumerate}[label=\alph*)]
  \item Prüfen Sie mithilfe der führenden Hauptminoren, dass $A$ positiv definit ist.
  \item Bestimmen Sie die Cholesky-Zerlegung $A = LL\transp$.
\end{enumerate}
\karo{7cm}

%==================================================================
\section*{Aufgabe 7 -- QR-Zerlegung\schwierig{mittel}}
Gegeben seien
\[
  Q = \frac{1}{5}\begin{pmatrix} 3 & -4 \\ 4 & 3 \end{pmatrix},\qquad
  R = \begin{pmatrix} 5 & 10 \\ 0 & 5 \end{pmatrix},\qquad
  A = QR,\qquad
  b = \begin{pmatrix} 5 \\ 15 \end{pmatrix}.
\]
\begin{enumerate}[label=\alph*)]
  \item Weisen Sie nach, dass $Q$ orthogonal ist, und bestimmen Sie $A = QR$.
  \item Lösen Sie das Gleichungssystem $A\vec{x} = b$ mithilfe der QR-Zerlegung.
  \item Bestimmen Sie $Q^{-1}$.
\end{enumerate}
\karo{7.5cm}

%==================================================================
\section*{Aufgabe 8 -- Eigenwerte und Diagonalisierung\schwierig{mittel}}
Gegeben sei
\[
  A = \begin{pmatrix} 4 & 1 \\ 2 & 3 \end{pmatrix}.
\]
\begin{enumerate}[label=\alph*)]
  \item Bestimmen Sie die Eigenwerte über das charakteristische Polynom $\chi_A(\lambda)$.
  \item Bestimmen Sie zu jedem Eigenwert einen Eigenvektor.
  \item Geben Sie $V$ und $D$ an mit $A = VDV^{-1}$ und berechnen Sie $V^{-1}$.
  \item Überprüfen Sie $\spur(A) = \sum_i \lambda_i$ und $\det(A) = \prod_i \lambda_i$.
\end{enumerate}
\karo{8cm}

%==================================================================
\section*{Aufgabe 9 -- Spektralzerlegung\schwierig{schwer}}
Gegeben sei die symmetrische Matrix
\[
  A = \begin{pmatrix} 2 & 1 \\ 1 & 2 \end{pmatrix}.
\]
\begin{enumerate}[label=\alph*)]
  \item Bestimmen Sie die Eigenwerte und eine \emph{orthonormale} Basis aus Eigenvektoren.
  \item Geben Sie die Spektralzerlegung $A = Q\Lambda Q\transp$ an
        (mit orthogonaler Matrix $Q$).
  \item Bestimmen Sie mithilfe der Spektralzerlegung die Inverse $A^{-1} = Q\Lambda^{-1}Q\transp$.
\end{enumerate}
\karo{7.5cm}

%==================================================================
\section*{Aufgabe 10 -- SVD und Singulärwerte\schwierig{schwer}}
Gegeben sei
\[
  A = \begin{pmatrix} 2 & 0 \\ 0 & 0 \\ 0 & 5 \end{pmatrix}.
\]
\begin{enumerate}[label=\alph*)]
  \item Berechnen Sie $A\transp A$ und bestimmen Sie daraus die Singulärwerte von $A$.
  \item Bestimmen Sie die Matrizen $V$, $\Sigma$ und $U$ der Singulärwertzerlegung
        $A = U\Sigma V\transp$.
  \item Bestätigen Sie durch Ausmultiplizieren, dass $U\Sigma V\transp = A$ gilt.
\end{enumerate}
\karo{8.5cm}

%==================================================================
\section*{Aufgabe 11 -- Pseudoinverse\schwierig{mittel}}
Gegeben sei die Matrix mit vollem Spaltenrang
\[
  A = \begin{pmatrix} 1 & 1 \\ 1 & 2 \\ 1 & 3 \end{pmatrix}.
\]
\begin{enumerate}[label=\alph*)]
  \item Berechnen Sie $A\transp A$ und $(A\transp A)^{-1}$.
  \item Bestimmen Sie die Moore-Penrose-Pseudoinverse $A^{+} = (A\transp A)^{-1}A\transp$.
  \item Verifizieren Sie, dass $A^{+}A = E$ gilt (Linksinverse).
\end{enumerate}
\karo{7.5cm}

%==================================================================
\section*{Aufgabe 12 -- Definitheit\schwierig{mittel}}
Gegeben sei
\[
  A = \begin{pmatrix} 2 & -1 & 0 \\ -1 & 2 & -1 \\ 0 & -1 & 2 \end{pmatrix}.
\]
\begin{enumerate}[label=\alph*)]
  \item Bestimmen Sie die führenden Hauptminoren $\Delta_1, \Delta_2, \Delta_3$ und
        entscheiden Sie damit über die Definitheit.
  \item Bestimmen Sie die Eigenwerte von $A$ und bestätigen Sie das Ergebnis aus a).
\end{enumerate}
\karo{7cm}

%==================================================================
\section*{Aufgabe 13 -- Hyperebenen\schwierig{mittel}}
Gegeben sei die Ebene
\[
  E:\quad 2x_1 - x_2 + 2x_3 = 3 .
\]
\begin{enumerate}[label=\alph*)]
  \item Geben Sie einen Normalenvektor $\vec{n}$ an und nennen Sie die Dimension von $E$ als
        Hyperebene im $\mathbb{R}^3$. Verläuft $E$ durch den Ursprung?
  \item Liegt der Punkt $p = (1,2,3)\transp$ auf $E$?
  \item Berechnen Sie den Abstand $d$ des Punktes $p$ zur Ebene $E$.
\end{enumerate}
\karo{6.5cm}

%==================================================================
\section*{Aufgabe 14 -- Finite Differenzen\schwierig{schwer}}
Gegeben sei das Randwertproblem
\[
  -u''(x) = 2 ,\qquad u(0) = u(1) = 0 ,
\]
mit den inneren Gitterpunkten $x_1 = 0{,}25,\; x_2 = 0{,}5,\; x_3 = 0{,}75$
(Schrittweite $h = 0{,}25$).
\begin{enumerate}[label=\alph*)]
  \item Leiten Sie mit dem zentralen Differenzenquotienten das lineare Gleichungssystem
        $A\vec{u} = \vec{b}$ her (geben Sie $A$ und $\vec{b}$ an).
  \item Nennen Sie vier strukturelle Eigenschaften der Matrix $A$.
  \item Lösen Sie das Gleichungssystem und bestimmen Sie $u_1, u_2, u_3$.
\end{enumerate}
\karo{8.5cm}

%==================================================================
\section*{Aufgabe 15 -- Lineare Regression\schwierig{schwer}}
Zu den Datenpunkten $(x,y) \in \{(1,1),\,(2,2),\,(3,2)\}$ soll eine Ausgleichsgerade
$y = c_0 + c_1 x$ bestimmt werden.
\begin{enumerate}[label=\alph*)]
  \item Stellen Sie die Designmatrix $X$ und den Vektor $\vec{y}$ auf und notieren Sie die
        Normalengleichung $X\transp X\,\vec{c} = X\transp \vec{y}$.
  \item Lösen Sie die Normalengleichung und bestimmen Sie $c_0$ und $c_1$
        (Hinweis: $\vec{c} = (X\transp X)^{-1} X\transp \vec{y}$).
  \item Geben Sie die Gleichung der Ausgleichsgeraden an.
\end{enumerate}
\karo{7.5cm}

\end{document}
